% Methods section — draft for review.
% Approach only: no counts, percentages, or outcome figures (those belong in
% Results). Class-agnostic: standard \section/\subsection and \cite{} keys.
% Citation keys are placeholders pending a .bib file:
%   openalex, genderize, santamaria2018, gbd2023, worldbank2024,
%   cohen1960, landis1977, duckdb, streamlit.

\section{Methods}

\subsection{Corpus definition}

The unit of analysis is the peer-reviewed research article, and the corpus is
defined by journal of publication. A journal was included if it met all of the
following: (i) global, international, or global public health is the primary
subject of its stated scope, which distinguishes dedicated global health
journals from general medical journals and from single-disease or
single-region specialty journals; (ii) it is peer-reviewed and indexed in a
major bibliographic database (MEDLINE, Scopus, or Web of Science); and (iii) it
published across the study window. The journals meeting these criteria are
listed with their ISSNs in Table~\ref{tab:journals}. Together they span the
orientations of the field: broad global health (\emph{Lancet Global Health},
\emph{BMJ Global Health}, \emph{Journal of Global Health}, \emph{Annals of
Global Health}, \emph{PLOS Global Public Health}); health policy, systems, and
economics (\emph{Health Policy and Planning}, \emph{Bulletin of the World
Health Organization}); implementation (\emph{Global Health Science and
Practice}); social determinants and health inequalities (\emph{Globalization
and Health}, \emph{Global Public Health}); and infectious disease in low- and
middle-income settings (\emph{Tropical Medicine \& International Health}).

A journal-defined boundary was chosen because it is transparent and
reproducible: the corpus is fully specified by the journal list and can be
re-retrieved exactly, whereas a keyword- or concept-based boundary depends on
query choices and on changes in indexing that make the denominator of every
subsequent proportion contestable. The trade-off is coverage: global health
research published in general medical journals (for example, \emph{The Lancet},
\emph{The New England Journal of Medicine}, and \emph{JAMA}) is not captured,
which is addressed in the Limitations.

The corpus comprises the works that OpenAlex~\cite{openalex} indexes in these
journals with a publication year from 2010 through the July 2026 retrieval
snapshot; 2026 is a partial year and is flagged in per-year analyses. Citation
counts and funding records reflect the state of OpenAlex at the snapshot date.

\subsection{Abstract processing}

OpenAlex stores abstracts as inverted indexes, which were reconstructed to
plain text. A work was treated as classifiable only if it carried a usable
abstract. Works with no abstract, with journal boilerplate in place of an
abstract, or with an abstract too short to classify were identified, recorded
by the reason their abstract was unusable, and excluded from classification
rather than treated as missing at random. Because unusable abstracts are not
distributed evenly across journals, analyses that depend on classification
under-represent the journals in which they concentrate; analyses that do not
depend on the abstract (publication counts, authorship, funding, institutions)
use the full corpus.

\subsection{Classification}

Each classifiable work was labeled over its title and abstract using a large
language model (Claude Opus 4.8, Anthropic), in successive passes: a topic
pass, a study-design (methods) pass, and a study-country pass. Each pass drew
on a defined taxonomy (Appendix): the topic pass assigns a primary category and
a finer subtopic and, where a second focus is co-equal, an optional secondary
category; the methods pass assigns a study-design type; and the study-country
pass assigns the country or countries in which the research was conducted, as
ISO 3166-1 alpha-2 codes, with a designated value for region-wide or global
work that names no specific country and another for cases that cannot be
determined. Each label carries a self-reported confidence level.

The prompts were generated from the taxonomy definition files, so that the
taxonomy is the single source of each label set. Model responses were parsed
and checked against the taxonomy; responses that were malformed or fell outside
the taxonomy were assigned the taxonomy's uncategorized value rather than
admitted as data. Study country denotes the location of the study population or
data collection, not author affiliation, and the prompt instructed the model
not to infer it from affiliations. The secondary topic category is a
descriptive enrichment layer, reported separately from the primary topic and
not used in the primary analyses. The full prompts and model settings are in
the repository.

\subsection{Validation}

Classification accuracy was assessed against a held-out sample of works
hand-labeled by a single human coder who did not see the model's labels, so
that the model's output could not bias the reference labels. The sample was
drawn from across the corpus journals, and works without a genuine abstract
were skipped. The labels evaluated were the classifier output as shipped; no
sampled work was used to develop the classifier, so the sample is a clean
held-out test of every label axis. Human and model labels were compared with
Cohen's $\kappa$~\cite{cohen1960}, interpreted using the Landis and Koch
descriptors~\cite{landis1977}; for study country, which may list more than one
country, exact-set agreement is reported alongside the coefficient. The
comparison is reproducible from the released labels and script.

\subsection{Author gender}

The likely gender of the first and last authors was inferred from their given
names using the Genderize.io service~\cite{genderize}, following two
adjustments recommended by Santamar\'ia and Mihaljevi\'c~\cite{santamaria2018}:
diacritics were removed, and for compound given names the first component was
used. Names returned below a set probability, single initials, and names for
which no gender could be determined were recorded as unknown. The inference is
probabilistic and binary and is less accurate for non-Western names; results
are reported only at the level of the corpus, never for individual authors.

\subsection{Funding}

Funding was taken from the grant records that OpenAlex stores for each work.
Each acknowledged funder was matched to a curated canonical funder list by its
OpenAlex funder identifier, with name aliases as a fallback, so that variant
names for the same funder were combined. The completeness of funder metadata
varies over time and is weaker in earlier years; it is therefore treated as a
measured quantity rather than assumed complete, and unfunded rates are read as
lower bounds.

\subsection{Disease burden}

Disease burden was taken from the Institute for Health Metrics and Evaluation
Global Burden of Disease 2023 results~\cite{gbd2023}, using
disability-adjusted life years and deaths (all ages, both sexes). Each topic
category was mapped to its corresponding GBD cause groups through a hand-built
crosswalk. Comparisons of research attention to disease burden use the pooled
World Bank low- and middle-income aggregate (the low, lower-middle, and
upper-middle income groups combined), because the corpus concerns global health
research directed largely at those settings. Burden values are fixed at the GBD
2023 release regardless of the corpus window.

\subsection{Analytic definitions and derived measures}

\paragraph{Analytic subset.}
Analyses of topics, methods, geography, funding, and institutions were
restricted to research articles: works classified as narrative reviews, as
commentary, editorials, or perspectives, or as undeterminable in design were
excluded, as were works with an uncategorized topic. The corpus overview
reports all works, and a separate analysis profiles the commentary and
editorial works on their own.

\paragraph{Leadership.}
An institution or country was credited with a work when one of its researchers
was the first or last author, the positions that conventionally denote study
leadership in health research. Middle-author positions were not counted toward
leadership.

\paragraph{Institution identity.}
Institutions were identified by their OpenAlex institution identifier rather
than by display name, because one display name (for example, ``Ministry of
Health'') is shared by many distinct national bodies. Where a name was shared
across multiple identifiers, the country was appended so that each national
entity is counted separately.

\paragraph{Income groups.}
Institutions and study countries were grouped by the World Bank income
classification for the 2024 fiscal year~\cite{worldbank2024}. Analyses contrast
high-income with combined low- and middle-income countries; where noted, the
World Bank tiers (low, lower-middle, upper-middle, and high income) are shown
separately.

\paragraph{Concentration.}
Concentration of funding and of institutional production was summarized with
the Gini coefficient and with the share held by the most productive
institutions.

\paragraph{Research leadership by geography.}
For research about a single country, leadership was classified from author
affiliations relative to that country: local-led (the last author is affiliated
with the study country), collaborative (an author, but not the last, is
affiliated with the study country), or externally led (no author is affiliated
with the study country). Studies spanning more than one country were excluded
from this classification.

\paragraph{Attention relative to burden.}
For each topic, its share of publications was compared with its share of low-
and middle-income disease burden. The two shares measure different quantities;
the comparison is descriptive and is not interpreted as a causal or normative
target.

\paragraph{Evidence strength.}
The study-design types were grouped into tiers for analysis (experimental and
synthesis designs, observational and analytic designs, and descriptive and
exploratory designs). The grouping is given in the taxonomy files.

\paragraph{Citation influence.}
Influence was measured with the OpenAlex cited-by count at the snapshot date.
Because citations accumulate over time, recent works are cited less than older
works of equal influence, and citation-based comparisons are read with that in
mind.

\paragraph{Collaboration.}
Collaboration breadth was measured as the number of distinct partner
institutions, identified by OpenAlex identifier, with which an institution
shared authorship on at least one work.

\subsection{Reproducibility and software}

The retrieval, classification, enrichment, and validation steps are implemented
as scripts; the taxonomies, the canonical funder list, the disease burden
crosswalk, the validation labels, and the agreement script are included in the
repository, and the corpus is rebuildable from OpenAlex. Analyses were
materialized into summary tables and are presented through an interactive
dashboard covering the corpus overview; the funding, geographic, topic,
methods, and institution analyses; and the abstract-availability report. Data
were held and queried in DuckDB~\cite{duckdb}, and the dashboard was built with
Streamlit~\cite{streamlit}. The code, data, and dashboard are available at
\url{https://github.com/ifeldhaus/global-health-research-map}.
